\begin{description}
	\item[dataブロック] データ長L，各時点の体重Wがデータになります。
	\item[parametersブロック] データ長と同じだけの状態を表す\verb|mu|と，状態の変動幅\verb|tau|，状態に付加される計測誤差の変動幅\verb|sig|を宣言しています。それに加えて，最初の状態$\mu_0$をパラメータとしました。これは2時点目の状態$\mu_2$は$\mu_2 \sim N(\mu_1,\tau)$で，3時点目は$\mu_3 \sim N(\mu_2,\tau)$で・・・と表現できるのに対し，最初の点は$\mu_1 \sim N(\mu_0, \tau)$となって，計測される前の状態を考えなければならないからです。わからないものは確率で表現する，というベイズの流儀に則って，これはパラメータとして推定してやることにします。
	\item[modelブロック] まず1時点目の状態は，0時点目の状態から出てくるものですので，$\mu_1 \sim N(\mu_0,\tau)$をモデル化しました。次に計測された値$W_l$は状態$\mu_l$から出てくるものですから，1時点目からデータ長Lまで順に$W_l \sim N(\mu_l,\tau)$として尤度を書きます。次に状態のモデルですが，\underline{2時点目以降}は前の時点からの影響で表現できるので，$\mu_i \sim N(\mu_{i-1}, \tau)$としています。最後に事前分布として，これまでの経験から$\mu_0 \sim N(80,10)$とし，変化の幅はいずれもSDですから半コーシー分布で表現しました。
\end{description}
